\textbf{Attribution map / heatmap} - A two-dimensional visualization
derived from a trained classifier that assigns relative importance to
image regions for a selected output score. In this thesis, attribution
maps are class-specific evidence maps for the selected pneumothorax or
hemorrhage target, not automatic anatomical segmentations.

\textbf{Positive evidence} - The part of a signed attribution map that
contributes positively toward the selected target output. For the CXR
pneumothorax experiments, positive evidence supports the model's
pneumothorax score.

\textbf{Negative evidence} - The part of a signed attribution map that
contributes negatively toward the selected target output. Negative
evidence is interpreted as suppressive evidence against the target score
and is considered significant for potential improvement when it overlaps 
the lesion mask.

\textbf{Magnitude view} - The absolute-value view of a signed
attribution map. It answers which pixels were influential, without
stating whether their influence increased or decreased the target score.

\textbf{Signed view} - A diverging visualization of the normalized
signed attribution map. It summarizes the local balance between positive
and negative evidence but does not replace inspection of the separate
positive, negative, and magnitude views.

\textbf{Selected view} - A visualization of other views using their absolute values and disregarding the alpha value.

\textbf{Localization validation} - Evaluation of whether
high-attribution pixels spatially correspond to available lesion
annotations, using metrics such as IoU, Dice, pointing-hit, and
precision-at-fraction.

\textbf{Faithfulness validation} - Evaluation of whether the highlighted
pixels materially affect the model's own output probability when they
are removed, replaced, or inserted. Faithfulness measures model
behavior, not clinical correctness.

\textbf{Top-fraction thresholding} - Conversion of a continuous heatmap
into a binary selected region by keeping a fixed fraction of the
highest-valued pixels. This makes heatmaps comparable to masks but
introduces a threshold choice that must be calibrated or reported.

\textbf{Coverage saturation} - A secondary diagnostic describing how
quickly top-fraction thresholding expands to cover nearly the whole
image. The proposed \texttt{coverage\_saturation\_fraction\_95} records
the smallest swept top-fraction at which selected coverage reaches at
least \texttt{95\%}; lower values indicate more diffuse or saturated
maps.

\textbf{Radiologist-centered review} - A structured qualitative
assessment in which a medically trained reader scores explanation
localization, usefulness, and failure mode categories under a fixed
review task. In this thesis it is an early controlled assessment, not
prospective clinical validation.
